\chapter{从受精卵到身体}

\begin{kaichang}
下次家里打生鸡蛋的时候，先别急着搅散。把蛋黄轻轻倒进一只平盘，仔细看蛋黄表面：十有八九你能找到一个小小的白点，直径只有两三毫米，像一粒粘在蛋黄上的米。养鸡人叫它“胚盘”。

如果这枚鸡蛋受过精，又被放在合适的温度里孵上二十一天，这个比芝麻大不了多少的白点，会长成一只完整的小鸡：有跳动的心脏、会啄壳的喙、两只眼睛、一身绒毛，破壳而出。没有经过任何“组装”，没有谁往里面塞零件，一只小鸡就从一个小点里“长”了出来。

先猜一猜：那个白点里的细胞，是怎么“知道”头该长在哪头、眼睛该长在哪、心脏该偏左一点的？答案写在心里，这一章结束时回来对——而谜底的核心，是你在前几章已经见过的东西：基因怎样被打开和关上。
\end{kaichang}

\section{白点里没有一个“小小鸡”}

先说一件人类猜错了两千多年的事。

在没有显微镜的时代，人们只能看到结果：蛋孵出鸟，种子长成树，孕妇的肚子一天天隆起。于是一个非常自然的猜想流传开来——也许那个白点（或者精子、卵子）里面，本来就藏着一只缩小版的完整小鸡，发育只是把它“吹大”而已。这套想法叫\emph{预成论}。十七世纪，甚至真的有显微镜学者声称在精子里看见了蜷缩的“小人”，还认真地画了下来。

如果预成论是对的，发育就没什么可研究的：一切早已画好，只等放大。但后来显微镜越来越好，人们盯着急剧变化的蛙卵和鸡胚一看，看到的不是小动物的放大，而是一连串\textbf{从无到有}的事件：一个圆球先变成一摞细胞，细胞球凹陷、折叠，长出沟槽和管子，某些地方隆起、某些地方断开——结构是在过程中一步步造出来的，之前根本不存在。这种观点叫\emph{渐成论}，它是对的。发育不是冲印照片，是现场施工。

那么施工现场在哪？就是开场那个白点。受过精的鸡蛋里，胚盘是受精卵（以及它分裂出的一小片细胞）；蛋黄只是随身的口粮。人的故事也一样：精子和卵子结合成受精卵，随后沿着输卵管一边旅行一边分裂，大约一周时着床，此后按照一张相当精确的时间表推进——

\begin{table}[htbp]
\centering
\begin{tabular}{ll}
\toprule
受精后时间 & 正在发生的大事 \\
\midrule
第 1 天 & 一个细胞，直径约 $0.1$ 毫米 \\
第 1 周 & 分裂成上百个细胞的小球，着床 \\
第 3--4 周 & 细胞球折叠出三层，神经管开始合拢 \\
第 5--8 周 & 心脏跳动，头尾分明，手指脚趾“雕刻”出来 \\
第 9 周以后 & 各器官长大、细化，体重从克级长到千克级 \\
\bottomrule
\end{tabular}
\end{table}

注意这张表的顺序：先有大致布局（头在哪、背在哪、四肢在哪），后有精雕细琢。先看整体、再抠细节，这正是理解发育的钥匙。

把镜头再推进一层，看看粒子尺度的施工现场。胚胎里的每个细胞，都是第 57 章认识的那座“化学城市”：细胞核里保存着同一份 DNA，细胞质里核糖体在翻译蛋白，细胞骨架的蛋白纤维在组装和拆解，细胞膜上的受体在监听外界。所谓发育，说到底就是几十万亿座这样的城市，在同一套化学规则下，各自做出不同的局部决定，最后拼出一个全局有序的身体。问题因此变得非常具体：这些局部决定——分裂还是休息、变成什么、搬去哪、活还是死——是被什么信息驱动的？

\section{一个细胞变成几十万亿个}

发育的第一道工序其实很“笨”：分裂，再分裂。有丝分裂你在第 64 章已经见过——染色体复制一套，分给两个子细胞，遗传信息原样传递。所以有个重要事实先摆在这：你身上几乎每个细胞，从脑细胞到脚趾甲细胞，携带的 DNA 序列基本相同。差异不在“原料清单”，而在“正在使用清单的哪几页”——第 70 章讲过的基因调控，和第 74 章讲过的表观遗传记忆，正是细胞各不相同的根源。

\begin{qiaoliang}
一个细胞要分裂多少次才够盖一个成年人？新生婴儿大约有 $2\times 10^{12}$（两万亿）个细胞，成年人约有 $3\times 10^{13}$ 到 $4\times 10^{13}$ 个。每次分裂细胞数翻倍，设分裂 $n$ 次，则 $2^n \approx 4\times 10^{13}$。两边取对数：$n \approx 13.6/0.3 \approx 45$。也就是说，\textbf{四五十轮分裂，就足以从一个细胞造出整个人体}。这个数字小得惊人——诀窍全在指数增长：第 10 轮才一千个细胞，第 20 轮一百万，第 30 轮十亿，第 40 轮一万亿。当然，真实情况不是整齐的同步分裂，不同部位分裂轮数差别很大，有些细胞（如多数神经元）很早就停止分裂，有些（如肠道上皮、骨髓）终身都在分裂更新（第 59 章讲过这种组织更新）。数量级估算的价值不在于精确，而在于让你看到：建造身体的“工程量”并没有想象中那么不可思议，难的不是数量，是\emph{排列}。
\end{qiaoliang}

这正是本章真正的问题：几十万亿个基因相同的细胞，凭什么有的变成神经元、有的变成肝细胞，还都长在了正确的位置上？

\section{地址从哪来：妈妈预先留下的坐标}

想象你蒙着眼睛被放进一座陌生城市，要找到自己的位置，你需要坐标。胚胎里的细胞也一样。而第一套坐标，是妈妈在卵子里提前放好的。

卵细胞看起来是个圆球，内部却并不均匀。以研究得最透的果蝇为例：妈妈的身体在制造卵细胞时，会把一些特定的信使 RNA（第 68 章讲过，这是基因的工作副本）和蛋白质\textbf{不对称地}堆放在卵的一端。有一种叫 Bicoid 的蛋白质，它的 RNA 被固定在卵的前端，受精后被翻译出来，蛋白质从前端向后端慢慢扩散，形成一条\textbf{前高后低的浓度斜坡}——像一滴墨水滴在水杯一角，还没搅匀时的样子。

这样的分子被叫做\emph{形态发生素}：它的浓度本身就是位置信息。胚胎前部的细胞“尝到”高浓度，中部尝到中等浓度，后部尝到低浓度。细胞不用看见全局，只要测量身边的浓度，就知道自己大概在身体的哪一段——这套办法有个形象的名字，叫\emph{法国国旗模型}：一条浓度渐变带，被两道“浓度门槛”切成蓝、白、红三块，每块的细胞执行不同的发育指令。

\begin{center}
\begin{tikzpicture}[scale=0.9]
\draw[thick] (0,0) rectangle (9,1.2);
\draw (3,0) -- (3,1.2);
\draw (6,0) -- (6,1.2);
\node at (1.5,0.6) {头部程序};
\node at (4.5,0.6) {躯干程序};
\node at (7.5,0.6) {尾部程序};
\draw[->] (0,-0.3) -- (9,-0.3) node[right]{位置（前 $\rightarrow$ 后）};
\draw[thick] (0,-1.5) .. controls (3,-1.2) and (6,-0.8) .. (9,-0.6);
\node at (4.5,-1.9) {形态发生素浓度：前高后低};
\draw[dashed] (3,-0.3) -- (3,0);
\draw[dashed] (6,-0.3) -- (6,0);
\end{tikzpicture}
\end{center}

（这是人为的表示法：真实的胚胎里既没有颜色，也没有直线边界，门槛两侧的细胞是在统计意义上倾向不同的命运。）

这套“浓度当坐标”的逻辑，解决了一个深刻难题：\textbf{全局图案可以只用局部规则生成}。没有一个“总指挥”细胞告诉每个细胞你是谁，只有扩散、浓度门槛和基因开关的连锁反应，秩序自己长了出来。

\section{基因网络的接力赛}

坐标只是开头。接下来是一整场精心编排的接力。

在果蝇胚胎里，母源的 Bicoid 蛋白是一种转录因子（第 70 章的角色：专门开关别的基因）。高浓度的 Bicoid 打开一批“大区间基因”，把胚胎分成几大段；这些基因的产物又是转录因子，再去开关下一批基因，把大段切成更细的条纹；条纹基因再切出体节的边界……一层激活一层，像涟漪扩散，又像裁判逐级细化赛道。最终，胚胎表面被划分成一连串重复的段落，每段细胞都拿到了自己的“段落编号”。

接力赛的下一棒，是一类大名鼎鼎的基因：\textbf{HOX 基因}。它们的工作是给每个段落发“身份牌”——这一段长触角，那一段长翅膀，再一段长平衡棒。HOX 基因全部编码转录因子，蛋白产物里都带一段约 60 个氨基酸、专门负责抓住 DNA 的结构（叫同源盒结构域），所以它们彼此长得很像，一看就是一大家子。

这一章第一次需要符号了，请记住两个要点：

\begin{itemize}[itemsep=2pt]
\item HOX 基因在染色体上\textbf{排成一列}，而且排列顺序和它们在身体上前 $\rightarrow$ 后表达的顺序一一对应：排最前的管最前部，排最后的管最后部。染色体上的“座位表”直接映射成身体的“楼层图”，这叫做\emph{共线性}，至今仍是发育生物学最漂亮的规律之一。
\item 这套基因古老得惊人。果蝇的 HOX 和人的 HOX（人有四簇，记作 \textit{HOXA}、\textit{HOXB}、\textit{HOXC}、\textit{HOXD}，分列在四条染色体上）是同一套祖传工具的变体。把人的某些 HOX 基因放进果蝇体内，居然能部分替代果蝇自己的版本工作——相隔六亿年进化的两个物种，身体布局用的还是同一套“地址系统”。
\end{itemize}

HOX 出错的后果极具戏剧性：果蝇的某突变让本该长平衡棒（一对小棍，飞行时用来感知旋转）的体节拿到了“长翅膀”的身份牌，于是长出两对翅膀；另一个突变让触角位置长出了腿。一个基因的开关错位，器官就“张冠李戴”——反过来说明，正常情况下器官位置就是被这类基因\textbf{指定}出来的。

\section{增殖、迁移、分化、凋亡：四种动作搭出身体}

有了坐标和身份牌，细胞实际干什么？归纳起来，发育只用了四种基本动作：

\begin{itemize}[itemsep=2pt]
\item \textbf{增殖}：有丝分裂，让细胞数量指数增长（前面算过，四五十轮就够）。哪些区域分裂快、哪些慢，直接决定器官的大小比例。
\item \textbf{分化}：同样的基因组，不同的使用方案。转录因子组合决定一个细胞开哪些基因、关哪些基因，最终变成神经元、肌肉细胞还是皮肤。第 74 章的表观遗传修饰，负责把这个决定“记”下来，让肝细胞的后代还是肝细胞。
\item \textbf{迁移}：细胞会“搬家”。脊椎动物胚胎里有一群叫神经嵴的细胞，从背侧的神经管出发，长途跋涉到全身各处，变成面部骨骼、色素细胞、部分神经——你的脸型，部分取决于几群细胞几周内的“徒步旅行”走没走对路。
\item \textbf{凋亡}：程序性细胞死亡。注意，这不是事故，是\textbf{施工手段}。人的手指就是这样“雕”出来的：胚胎早期手是一片连着的“桨”，指间细胞按计划集体凋亡，五指才分开。凋亡程序失灵一点点，结果就是并指。大脑发育中也有大量神经元按程序退场——先过量生产，再按连接质量淘汰。
\end{itemize}

四种动作全都受基因网络调度，又反过来改变细胞的位置和邻居，从而改变下一轮信号。发育就是这样一圈圈推进的动态过程，任何一张静态“图纸”都装不下它。

生活里最直观的证据藏在伤口里。你擦破皮后，伤口边缘的细胞接到信号，开始增殖（补上缺口）、迁移（向中间爬行）、分化（变成新的皮肤细胞），多余的细胞再被清理——四种动作原样重演一遍，只是规模小得多。发育从未真正“结束”，它只是从轰轰烈烈的施工期，转入了低强度的维护期（第 59、60 章讲过这种终身更新与稳态）。

\section{不只是基因：力学与环境也在说话}

到这一步，你可能会以为形态完全写在基因里。但还有两只手在塑形。

第一只是\textbf{力学}。细胞不是飘在汤里的独立颗粒，它们彼此黏连、推拉、挤压。早期胚胎会经历一场惊人的变形：一片细胞向内凹陷、卷入，像把气球的一角摁进去，形成原肠，身体的“内外”就此分开。驱动力来自细胞骨架（第 57 章）的收缩和细胞间黏附的变化——纯粹的物理力。反过来，细胞还能“感觉到”拉扯：受压的组织会改变基因表达。骨骼就是这种循环的产物——哪里常受压力，哪里的骨就长粗。基因给力学搭台，力学给基因反馈。

第二只是\textbf{环境}。温度、营养、化学物质都能插手发育。鳄鱼和海龟的性别不由染色体决定，而由孵化温度决定——同一窝蛋，差两三摄氏度，孵出的性别比例就不同。一个沉重的医学教训来自上世纪中叶：一种叫沙利度胺的止吐药被广泛用于缓解孕吐，结果导致上万名婴儿四肢发育严重受损——药物恰好干扰了四肢快速成型的那几周。这件事直接催生了现代孕期药物审查制度。它也说明一个规律：\textbf{发育中的身体对干扰极其敏感，而且敏感期有严格的时间窗口}——同样的干扰，早两周或晚两周，后果完全不同。这就是为什么医生对孕期用药、营养（比如叶酸补充预防神经管闭合缺陷）如此较真：剂量、时机、个体差异，每一项都可能改写施工进程。

\begin{shiyan}
\textbf{观察：找一找鸡蛋的胚盘}（安全，可在家做）

材料：生鸡蛋 1～2 枚、平盘一只、手电筒（可选）。

步骤：① 把鸡蛋轻轻磕开，让蛋黄完整滑入平盘，不要戳破。② 在蛋黄表面寻找那个直径两三毫米的灰白色小圆点——胚盘。③ 关掉顶灯，用手电筒从侧面照蛋黄，胚盘会显得更清楚。④ 对比一下：蛋清里那两条白色的螺旋状“带子”（系带）是固定蛋黄用的，不是胚胎。

观察点：胚盘相对于蛋黄有多小？超市里买的鸡蛋大多未受精，胚盘只是一个小白点；如果能找到受精蛋（如农家土鸡蛋），白点会稍大、边缘更规整。想一想：这么小的一点，要靠哪些信息才能长成一只小鸡的完整布局？

安全提示：生鸡蛋可能带沙门氏菌，做完后洗手，蛋液彻底做熟再吃或丢弃，不要生食。
\end{shiyan}

\section{科学家怎样知道：把果蝇翻个底朝天}

\begin{zhengju}
上世纪七十年代末，发育生物学面对一个尴尬局面：胚胎怎么长出花纹和体节，谁都只有模糊的猜想。有人猜是化学梯度，有人猜细胞按时间计数，各说各话——因为没人知道到底有哪些基因参与。

在德国海德堡的一间实验室里，遗传学家妮斯莱因–福尔哈德和维绍斯决定用最“笨”也最彻底的办法：\textbf{饱和筛选}。思路是这样的——如果某个基因对发育必不可少，把它破坏掉，胚胎就会在某个环节死掉，而且死相（花纹和体节的缺陷模式）会暴露这个基因原来的职责。于是他们给雄果蝇喂诱变剂（化学诱变会随机破坏 DNA，第 72 章讲过突变的来源），再一代代杂交，最后收集那些\textbf{纯合突变、在胚胎期死亡}的果蝇胚胎，一只只放到显微镜下看皮肤花纹。

这是一项体力活：他们筛查了几千个突变品系、数万只胚胎。但回报巨大。他们找到了一批“体节极性基因”：有的突变胚胎体节数减半，有的每一段花纹镜像重复，有的整条条纹消失——缺陷模式和正常花纹之间，呈现出清晰的对应关系。换句话说，胚胎的花纹不是混沌长成的，而是被一组有限、可列举的基因\textbf{分层控制}的。前面讲的“母源梯度—大区间—条纹—体节”接力赛，正是从这批突变体的死相里倒推出来的。

替代解释也被逐一排除。比如“这些缺陷只是胚胎普遍中毒的假象”——如果那样，所有突变体应该死得差不多，而不是各有各的、可重复的特征花纹。又如“梯度只是想象”——后来科学家真的做出了 Bicoid 蛋白的抗体，染出一条前高后低的真实梯度；再人为改变卵里 Bicoid 的浓度，胚胎的头部结构随之扩大或缩小，预言兑现。1995 年，这项工作获得诺贝尔生理学或医学奖。

故事还有下半场。刘易斯多年研究的 HOX 基因突变（两对翅膀的果蝇）与这批条纹基因接上了头：先分段，再给每段发身份牌。而当分子生物学技术成熟后，人们在老鼠、人甚至海胆体内都找到了同一套基因家族——果蝇身上的“手术刀”，划开的是整个动物界共用的发育工具箱。这正是本章最该记住的方法论：\textbf{想理解一台机器，最锋利的办法是逐个弄坏它的零件，看机器怎么坏}。
\end{zhengju}

\section{这套解释在哪里会失效}

本章的模型很好用，但请记住它的边界。

第一，\textbf{果蝇是特例，不是标准答案}。果蝇早期胚胎是一团共享细胞质的“合胞体”：细胞核分裂了很多轮，却不急着分隔成独立细胞，形态发生素可以在共同的细胞质里自由扩散，梯度特别好建立。人和小鼠的胚胎从第一次分裂起就是独立细胞，位置信息更多靠细胞之间传递信号（一个细胞释放、邻居接收），机制和果蝇形似而神不全同。直接从果蝇外推到人，会栽跟头。

第二，\textbf{“浓度门槛”是简化}。真实胚胎里，细胞读的不是单一形态的浓度，而是多种信号的组合、信号持续的时间、甚至信号变化的节奏。法国国旗模型抓住了“位置决定命运”的精髓，却装不下全部细节。

第三，\textbf{知道了基因网络，也预测不了每一根指纹}。发育里有真实的随机性：基因表达的分子涨落（第 70 章提过单细胞间的随机波动）、细胞迁移路径的微小差异，都会积累放大。同卵双胞胎基因完全相同，指纹却不同，患某些疾病的运气也不同——这正是红线之一：DNA 不是完整决定人的程序代码。本章解释的是“身体怎样可靠地长出大致布局”，不是“每个细节都被写死”。

\begin{wuqu}
“基因就是身体的建筑图纸，发育只是照图施工。”——图纸比喻流传极广，但它会误导你三件大事。其一，图纸在施工前就完整存在，而胚胎里没有任何地方存着一幅“全身图”：细胞只掌握局部规则（测浓度、听邻居），全局图案是涌现出来的。其二，图纸可以撕成两半各盖半栋楼，而 1891 年杜里舒把两细胞期的海胆胚胎一分为二，每一半都长成了完整（只是较小）的幼体——每个细胞带的是“全部规则”，不是“半张图”。其三，图纸盖楼是单向执行，发育却是双向对话：力学、营养、温度都会反过来修改基因的表达方案。更好的比喻是：基因组像一本菜谱合集加一套厨房规则，做出什么菜，取决于此刻哪些原料在手、厨师（细胞）处于什么位置、火候如何——同一套规则，位置不同，成品不同。
\end{wuqu}

\section{再看那个小白点}

现在回到开场那枚鸡蛋。胚盘里的细胞没有任何一个“看见”过整只鸡。它们拥有的只是：妈妈产卵时预先布置的不对称分子、受精后建立的前后与背腹坐标、一波波扩散的信号分子、层层叠叠的转录因子接力、HOX 家族发放的体段身份牌，以及四种动作——分裂、分化、搬家、按程序退场——在力学与环境的推拉中反复执行。

头长在前端，是因为前端的细胞读到了“前”的坐标，打开了头部程序；心脏偏左，是因为早期胚胎有一层细胞把一种信号分子只送到了左侧（这个“左右对称破缺”的机制本身，至今仍有细节在争论）；两只眼睛位置对称，是因为两侧细胞读到了相同的坐标。没有指挥，没有图纸，只有局部规则加上亿万年调试好的接力顺序——一只小鸡就这样从白点里长出来。

你也是一样。你曾是一个直径约 $0.1$ 毫米的细胞，经过四五十轮分裂、无数次开关切换和几场精心安排的细胞凋亡，长成了正在读这本书的身体。你的肚脐，就是那场施工留下的纪念章。

\section{往前一步：工具箱的来历}

理解发育，正在改变医学。科学家已经能在培养皿里用干细胞（第 59 章讲过它们的多能性）人为重现部分发育信号，养出“类器官”——米粒大小、却能部分模拟肠道、视网膜甚至大脑早期结构的细胞团，用来研究疾病、测试药物，而不必惊动真正的胚胎。先天畸形的预防（如孕期补叶酸）、对发育敏感期用药的严格管理，也全都建立在本章这套图景之上。

但请注意一个一直没解释的巨大事实：这套发育工具箱\textbf{古老得离谱}。HOX 基因在果蝇和你体内是同一套；制造眼睛的关键开关基因，在苍蝇、老鼠和人身上几乎通用——把老鼠的对应基因在果蝇翅膀上错误打开，翅膀上会长出一只果蝇式的眼睛。建造身体的基本规则，在几亿年的分异中几乎没有改写，只是被重新组合、重新排时间。

这就留下了一个无法回避的问题：这套工具箱最初是怎么来的？基因序列一代代传递时会发生变化（第 72 章的突变），有些变化改变了发育的时间和位置，于是一种动物慢慢变成另一种动物。谁在“筛选”这些变化？变化又怎样积累成物种之间的鸿沟？那是第十篇、也是下一章——第 76 章《自然选择不是“努力就会进化”》——要展开的故事。发育是基因在一代之内的演出；进化，是这套剧本在百万代里的修改史。

\begin{tiaozhan}
形态发生素梯度能解释“前高后低”的坐标，但它解释不了另一种常见图案：斑马条纹、豹子斑点、手指的五根并列——这些\textbf{周期性重复}的结构，单靠一条单调斜坡切不出来。1952 年，数学家图灵（对，就是计算机科学的那位图灵）提出一个惊人的方案：想象两种化学物质，一种叫激活剂，能促进自己也促进另一种抑制剂；抑制剂只负责压制激活剂，但扩散得更快。于是任何一处激活剂的偶然升高，都会自我加强形成一个“峰”，而它产生的抑制剂跑得更快，在周围压出一圈“谷”——峰与峰自动保持固定间距，斑点或条纹凭空出现，无需任何预设坐标。这套“反应–扩散”机制的数学只用两条规则：自我激活，加上“抑制跑得比激活快”。鱼身上的条纹会随鱼的长大而变宽，斑马鱼条纹被剪断后会按预测方式重连——半个多世纪后，图灵的纯数学猜想被实验一点点坐实。细节涉及微分方程，跳过不影响主线；但请记住这个思想：\textbf{图案不需要蓝图，只需要正确的局部规则}。
\end{tiaozhan}

\section*{练一练}

\begin{sikaoti}
\item 画一张信息流图：从“妈妈在卵前端放入 Bicoid 的 RNA”开始，到“胚胎某段细胞获得体节身份”结束，至少画出四个环节，并在每个箭头旁注明发生的是扩散、转录、翻译还是基因开关。
\item 用正文的方法估算：如果某器官从 1000 个细胞开始，需要长到约 100 万个细胞，大约还要分裂多少轮？（$2^{10} \approx 10^{3}$，可以直接用心算。）
\item 杜里舒把两细胞期的海胆胚胎分开，两半各自长成完整个体。如果用“图纸”比喻，预期结果应该是什么？这个实验为什么能推翻预成论，却不能推翻“基因很重要”？
\item 一位孕妇在不知怀孕的第 4 周和第 20 周分别接触了同一种发育干扰物。参照正文的时间表，预测哪一次风险更大，并解释“敏感期窗口”的含义。
\item 人的手指靠细胞凋亡“雕刻”出来。请再想出身体发育中两个可能用到“先多造、后淘汰”策略的例子（提示：一个在神经系统，一个在免疫系统，可回顾第 61 章），并说明这种浪费为什么可能是合理的。
\item 画一幅草图表示 HOX 基因的“共线性”：横轴是染色体上的基因排列，纵轴是身体前后位置，标出基因座位顺序与表达位置的对应关系，并在图旁注明这是一种表示法。
\end{sikaoti}
